\chapter{Conclusão}
\label{cap:conclusao}

Este trabalho partiu de um problema concreto observado na trajetória de estudantes de Ciência da Computação da UFSC: a distância entre a programação em Python, já familiar da disciplina de Programação Orientada a Objetos, e a descrição de \textit{hardware} em VHDL, exigida logo nos primeiros contatos com a disciplina de Circuitos e Técnicas Digitais. Diante disso, o objetivo geral proposto no Capítulo~\ref{cap:introducao} foi desenvolver, em Python, um simulador didático de circuitos digitais baseado em propagação reativa de sinais, capaz de apoiar o aprendizado dos conceitos introdutórios dessa disciplina.

Os cinco objetivos específicos definidos no Capítulo~\ref{cap:introducao} foram atingidos: (i) o barramento de \textit{bits} com notificação automática foi implementado pela classe \texttt{BitArray}; (ii) as portas lógicas básicas (\texttt{Not}, \texttt{And}, \texttt{Or}, \texttt{Xor}) foram implementadas como componentes reativos; (iii) o componente sequencial \texttt{BehavioralRegister}, sensível a borda de subida de \textit{clock} e escrito em lógica Python, foi implementado e validado isoladamente, assim como sua contraparte construída inteiramente com \textit{latches} (\texttt{SRLatch}, \texttt{DLatch}, \texttt{DFlipFlop} e sua extensão a N \textit{bits}, \texttt{DFlipFlopRegister}); (iv) a composição desses componentes em circuitos maiores foi demonstrada por meio de um meio-somador, um somador completo, três contadores binários (com o atalho \texttt{BehavioralAdder}, com \texttt{RippleCarryAdder} e \texttt{BehavioralRegister}, e com \texttt{RippleCarryAdder} e \texttt{DFlipFlopRegister}) e um \textit{flip-flop} \textit{toggle}; e (v) todos esses circuitos foram validados no Capítulo~\ref{cap:resultados}, inclusive por meio de testes automatizados, com a saída do simulador correspondendo exatamente ao comportamento esperado em cada caso, incluindo o tratamento de \textit{overflow} nas suas duas formas (truncamento no \texttt{BehavioralAdder} e sinal \texttt{carry\_out} no \texttt{RippleCarryAdder}).

De forma mais ampla, os resultados obtidos sustentam a hipótese levantada na motivação deste trabalho: um simulador implementado inteiramente em Python, sem dependências externas, é capaz de reproduzir corretamente os conceitos introdutórios de circuitos digitais em um ambiente de execução já familiar ao estudante, servindo como uma etapa intermediária de aprendizado antes do contato com HDLs tradicionais como VHDL e Verilog.

\section{Contribuições}

As principais contribuições deste trabalho são:

\begin{itemize}
    \item um simulador de circuitos digitais funcional, implementado em Python 3 sem dependências externas, organizado em quatro pacotes (\texttt{signals}, \texttt{gates}, \texttt{components}, \texttt{utils}) que separam claramente representação de sinais, portas lógicas elementares, componentes compostos e funções auxiliares;
    \item uma arquitetura reativa baseada no padrão de projeto \textit{Observer}, na qual cada componente atua simultaneamente como observador de suas entradas e sujeito observado por sua saída, propagando mudanças automaticamente ao longo do circuito;
    \item dois somadores com propósitos didáticos distintos e explicitamente documentados: o \texttt{RippleCarryAdder}, montado inteiramente a partir de portas lógicas e cujo \textit{overflow} surge naturalmente da topologia do circuito, e o \texttt{BehavioralAdder}, que demonstra a possibilidade de misturar Python ``puro'' com o restante do circuito reativo quando o objeto de estudo é outra parte do circuito;
    \item um registrador construído inteiramente com \textit{latches} (\texttt{SRLatch}, \texttt{DLatch}, \texttt{DFlipFlop} mestre-escravo, e sua extensão a N \textit{bits}, \texttt{DFlipFlopRegister}), complementando o \texttt{BehavioralRegister} e validado, inclusive, na configuração mais exigente testada neste trabalho (um \textit{flip-flop} \textit{toggle} realimentado através de si mesmo) e em um contador de 3 \textit{bits} completo, comparado lado a lado com o contador equivalente montado sobre o \texttt{BehavioralRegister};
    \item um conjunto de circuitos de exemplo (meio-somador, somador completo, três contadores binários e um \textit{flip-flop} \textit{toggle}, este último nas versões com \texttt{BehavioralRegister} e com \texttt{DFlipFlop} de portas) validados comparando a saída do simulador ao comportamento teoricamente esperado e verificados por um conjunto de testes automatizados;
    \item um conjunto de recursos de observação e depuração construídos sobre o próprio mecanismo de \textit{listeners} (visualização de estados transitórios, medição de atividade de chaveamento, sondas com asserções, reconstrução do caminho da propagação e verificação de equivalência entre descrições), escritos em poucas linhas de Python e comparados aos recursos equivalentes de outras ferramentas (Seção~\ref{sec:observacao}).
\end{itemize}

\section{Trabalhos Futuros}

Como possibilidades de continuidade para este trabalho, destaca-se a inclusão de novos componentes combinacionais, como multiplexadores e decodificadores, além de um sinal de \textit{reset} assíncrono para o \texttt{BehavioralRegister}. No plano da \textit{engine} de simulação, seria relevante detectar laços combinacionais inválidos automaticamente e introduzir um modelo de atraso de propagação (ciclos delta, ou um agendador de eventos com tempo simulado), tornando a estabilização de circuitos com realimentação através de \textit{latches} independente de uma convenção específica de ordenação de \textit{listeners} (Seção~\ref{sec:contador-portas-real}). Do ponto de vista de usabilidade, a geração de diagramas de tempo (por exemplo, exportando o histórico dos sinais em formato VCD para visualização em ferramentas como o GTKWave), um gerador de \textit{clock} automático (independente da alternância manual de \texttt{clock.output}) e uma interface gráfica/visual para acompanhar a propagação dos sinais em tempo real tornariam o simulador mais acessível a estudantes iniciantes. Por fim, como validação pedagógica mais ampla, propõe-se avaliar o uso do simulador em turmas reais da disciplina de Circuitos e Técnicas Digitais, comparando a curva de aprendizado com a abordagem tradicional baseada diretamente em VHDL.
